\documentclass[FGBook.tex]{subfiles}

\begin{document}

\drama{Romeo a Julie}{WILLIAM SHAKESPEARE}

\noindent\begin{wrapfigure}{r}{0.25\textwidth}\footnotesize
\subwection{Ukázka z díla:}
\setlength{\parindent}{3pt}
\noindent \textit{Jul.} Ach, žel! \\
\textit{Rom.} Teď mluví! -- Promluv poznovu, \par
ó jasný anděle! -- neb nad hlavou \par
mi vystupuješ zářná z noční tmy, \par
jak okřídlený posel nebeský \par
v sloup obráceným očím žasnoucích \par
smrtelných lidí, padajících na znak, \par
by dívali se naň, jak osedlal \par
si lenivě plynoucí oblaky \par
a po nebeských ňadrech vesluje. \\
\textit{Jul.} Romeo, Ó Romeo! -- proč's Romeo? \par
Své jméno zapři, otce zřekni se, \par
neb, nechceš-li, mně lásku přísahej, \par
a nechci dál být Capuletova.
\end{wrapfigure}

\wection{LITERÁRNÍ TEORIE}

\subwection{Literární druh a žánr:}
\noindent drama, tragédie (tragédie osudu s prvky romantické lásky a pomsty)

\subwection{Literární směr:}
\noindent Anglická renesance (elizabethanská doba), humanismus. Shakespeare kombinuje prvky středověké tragédie, petrarkovské lásky a italské novely.

\subwection{Jazyk a slovní zásoba:}
\noindent Blankvers (pětistopý jambický verš) u vyšších postav, proza u nižších vrstev (chůva). Bohatý poetický jazyk s metaforami, přirovnáními, slovními hříčkami a citově zabarvenými výrazy. Kontrast vysokého stylu (Romeo, Julie) a lidové mluvy.

\subwection{Figury:}
\noindent Řečnické otázky (\enquote{Romeo, proč's Romeo?}), metafora, přirovnání, oxymóron (sladká bolest lásky), dramatická ironie.

\subwection{Tropy:}
\noindent Symbolismus (balkon = hranice mezi rody, jed = osud), alegorie (feud = nenávist společnosti), foreshadowing (prolog sonet předpovídá tragický konec).

\subwection{Stylistická charakteristika textu:}
\noindent Pět dějství s chronologickým postupem. Směs tragického a komického (první akty mají komediální prvky). Rychlé dialogy, monology (balkonová scéna) a sbor (prolog). Tragikomický tón – láska a smrt jsou propojeny.

\subwection{Postavy:}
\noindent 
\textbfhl{ROMEO MONTEK --} mladý, romantický, impulzivní, zamilovaný do lásky samotné (nejprve Rosalina, pak Julie). Statečný, ale ukvapený.\\
\textbfhl{JULIE KAPULETOVÁ --} mladá (téměř 14 let), inteligentní, odvážná, tvrdohlavá. Zpočátku poslušná, později se vzepře rodičům kvůli lásce.\\
\textbfhl{TYBALT --} Juliin bratranec, horkokrevný, bojovný, symbol rodinné nenávisti.\\
\textbfhl{MERKUCIO --} Romeův přítel, vtipný, cynický, bojovný. Jeho smrt urychlí tragédii.\\
\textbfhl{BRATR VAVŘINEC --} františkán, moudrý, ale naivní. Snaží se pomoci, jeho plán selže.\\
\textbfhl{CHŮVA --} Juliině věrná služebnice, komická, lidová postava, praktická.

\subwection{Děj:}
\noindent Ve Veroně žijí dva znepřátelené rody – Montekové a Kapuletové. Romeo (Montek) se na plese zamiluje do Julie (Kapuletová). Oba se tajně oddají u bratra Vavřince. Po pouliční bitce Romeo zabije Tybalta (Juliina bratrance) a je vyhoštěn. Julie má být provdána za Parise. Bratr Vavřinec jí dá uspávací nápoj, aby vypadala mrtvá. Plán selže – Romeo se dozví o „smrti“ Julie, vrátí se, zabije Parise a spáchá sebevraždu jedem. Julie se probudí a probodne se Romeovou dýkou. Obě rodiny se nad mrtvými těly usmíří.

\subwection{Kompozice, prostor a čas:}
\noindent Pět dějství s prologem (sonet, který shrnuje děj). Chronologický postup (události proběhnou během několika dní). Prostor: Verona (Itálie). Čas: 14.–16. století (elizabethanská doba).

\subwection{Význam sdělení:}
\noindent Tragédie nešťastné lásky, kterou ničí rodinná nenávist. Síla lásky překonává společenské bariéry, ale osud („star-crossed lovers“) a ukvapenost vedou k tragédii. Kritika slepé nenávisti a rodičovské autority. Láska jako osvobozující, ale destruktivní síla. Závěrečné usmíření ukazuje, že smrt mladých přináší mír.

\vspace*{2em}\\
\wection{LITERÁRNÍ HISTORIE}

\subwection{Politická situace:}
\noindent Alžbětinská Anglie (vláda Alžběty I., konec 16. století) – náboženské napětí (katolíci vs. protestanti), stabilita po občanských válkách. Shakespeare kritizuje feudální spory skrz italské prostředí.

\subwection{Kontext dalších druhů umění:}
\noindent Inspirováno italskými novelami (Matteo Bandello) a anglickou básní Arthura Brooka (1562). Popularizovalo téma zakázané lásky.

\subwection{Kontext literárního vývoje:}
\noindent Raná Shakespearova tragédie. Kombinuje komediální prvky (první akty) s tragédií. Ovlivnil pozdější romantismus a moderní adaptace (West Side Story, filmy).

\subwection{Autor {\ssmall -- život autora:}}
\noindent William Shakespeare (1564 Stratford-upon-Avon – 1616). Nejvýznamnější anglický dramatik. Psal v alžbětinské době. Autor 37 her, 154 sonetů. Jeho tvorba se dělí na rané (komedie, historické hry), střední (tragédie jako Hamlet) a pozdní období.

\subwection{Další autorova tvorba:}
\noindent Mezi nejznámější díla patří \textit{Hamlet}, \textit{Othello}, \textit{Sen noci svatojánské}, \textit{Král Lear}, \textit{Makbeth}.

\subwection{Aktuálnost tématu a zpracování tématu:}
\noindent Nadčasové – zakázaná láska, generační konflikt, nenávist mezi skupinami. Adaptace do moderních kontextů (rasismus, gangy) ukazují, že téma zůstává relevantní.

\vspace*{2em}\\
\wection{OSOBNÍ NÁZOR}
\noindent Romeo a Julie je klasika, která stále zasáhne. Balkonová scéna a intenzita mladé lásky jsou nádherné, ale nejsilnější je tragédie způsobená rodinnou nenávistí a ukvapeností. Shakespeare mistrně ukazuje, jak láska může být osvobozující i destruktivní. Prolog sonet hned prozradí konec, přesto napětí funguje. Je to varování před slepou nenávistí a oslavou lásky, která překonává společenské bariéry. I po stovkách let patří k nejsilnějším milostným příběhům – doporučuji číst i zhlédnout.

\vfill

\noindent\begin{minipage}{\textwidth}
    {\textcolor{\wpagecolor}{\rule{\linewidth}{0.4pt}}
    \footnotesize
    \textbfhl{Tragédie --} dramatický žánr s tragickým koncem hrdinů a morálním ponaučením; 
    \textbfhl{Blankvers --} nerýmovaný pětistopý jambický verš; 
    \textbfhl{Foreshadowing --} předznamenání budoucích událostí (prolog).
    }
\end{minipage}

\newpage

\end{document}